\section*{ESPECIFICACIÓN DE REQUERIMIENTOS DE\\SOFTWARE (ERS)}
\begin{center}
\textit{Documento Técnico de Diseño y Contrato de Desarrollo}
\end{center}

\section{Control de Versiones}

\begin{tabularx}{\textwidth}{p{0.14\textwidth}p{0.18\textwidth}Xp{0.20\textwidth}}
\toprule
\rowcolor{azulprincipal}\textcolor{white}{Versión} & \textcolor{white}{Fecha} & \textcolor{white}{Descripción del Cambio} & \textcolor{white}{Autor}\\
\midrule
0.1 & 31/07/2026 & Línea base a partir de la sesión de captura SES-001. & Analista de Sistemas\\
\bottomrule
\end{tabularx}

\section{Requerimientos Funcionales (RF)}

\small
\begin{longtable}{>{\bfseries}p{0.11\textwidth}p{0.18\textwidth}p{0.34\textwidth}p{0.22\textwidth}}
\toprule
\rowcolor{azulprincipal}\textcolor{white}{ID} & \textcolor{white}{Nombre} & \textcolor{white}{Descripción Técnica Rigurosa} & \textcolor{white}{Criterio de Aceptación}\\
\midrule
RF-001 & Agendar citas & El sistema debe permitir agendar una cita seleccionando médico, paciente, fecha, hora, servicio y motivo de consulta. La cita debe quedar asociada a un estado inicial de ``Programada''. & Se registra la cita con un identificador único y se puede consultar en la agenda del médico y de recepción.\\
RF-002 & Validación de citas & El sistema debe comprobar que no exista otra cita activa para el mismo médico en la misma fecha y hora. La validación debe ejecutarse antes de guardar la cita. & Si existe un conflicto, la operación se rechaza y se informa al usuario que el médico ya tiene una cita asignada en ese intervalo.\\
RF-003 & Módulo de cobro & El sistema debe permitir registrar el cobro de una cita, incluyendo monto, fecha, medio de pago y estado de pago. & Se muestra el cobro asociado a la cita, con su monto y estado; además, puede incluirse en el arqueo diario.\\
RF-004 & Control de insumos & El sistema debe permitir registrar, consultar, modificar y gestionar insumos, incluyendo existencias actuales, stock mínimo, entradas, salidas y consumos clínicos. & Se muestra el stock actualizado de cada insumo y se genera una alerta cuando alcanza el stock mínimo configurado.\\
RF-005 & Control de medicamentos & El sistema debe permitir consultar los medicamentos sujetos a receta y vincular cada dispensación o venta con el paciente, la consulta y la prescripción médica correspondiente. & Se muestran los medicamentos y recetarios asociados a cada paciente; no se permite registrar la dispensación sin una prescripción válida.\\
\bottomrule
\end{longtable}
\normalsize

\section{Requerimientos No Funcionales (RNF)}

\small
\begin{longtable}{>{\bfseries}p{0.11\textwidth}p{0.17\textwidth}p{0.35\textwidth}p{0.22\textwidth}}
\toprule
\rowcolor{azulprincipal}\textcolor{white}{ID} & \textcolor{white}{Categoría} & \textcolor{white}{Especificación de Atributo} & \textcolor{white}{Métrica de Verificación}\\
\midrule
RNF-001 & Seguridad & La modificación de parámetros críticos de stock, incluido \texttt{stock\_minimo}, queda restringida a los roles Administrador o Gerente. & Prueba de control de acceso: un usuario con rol no autorizado recibe acceso denegado y la operación no modifica los datos.\\
RNF-002 & Rendimiento y tiempo de respuesta & Las consultas de agenda, pacientes e inventario deben responder en un máximo de 2 segundos bajo la carga esperada de al menos 200 pacientes semanales. & Prueba de rendimiento con datos representativos: al menos el 95\% de las consultas debe responder en 2 segundos o menos.\\
RNF-003 & Disponibilidad & El sistema debe estar disponible durante el horario de atención y contar con un procedimiento de recuperación ante fallos. & La disponibilidad objetivo será de al menos 99\% durante el horario operativo, verificada mediante registros de monitoreo.\\
RNF-004 & Sistema intuitivo & La interfaz debe ser comprensible para recepción, médicos y administración, incluyendo usuarios con poca experiencia informática. & Prueba de usabilidad con representantes de cada rol: deben completar las tareas principales sin asistencia crítica.\\
RNF-005 & Confidencialidad & Los datos personales y clínicos de los pacientes deben protegerse mediante autenticación, autorización por roles y cifrado durante la transmisión. & Revisión de permisos y pruebas de acceso indebido: un usuario no autorizado no puede consultar ni exportar información clínica.\\
\bottomrule
\end{longtable}
\normalsize

\section{Reglas de Negocio y Restricciones}

\subsection*{RN-01: Exclusividad de Agenda}
\textbf{Regla:} Un médico no puede tener asignada más de una cita médica dentro del mismo intervalo de tiempo. Para la agenda inicial se utilizarán bloques estándar de 30 minutos.

\textbf{Propósito:} Eliminar el problema de asignar dos pacientes a la misma hora al mismo médico.

\subsection*{RN-02: Control Inviolable de Stock Mínimo}
\textbf{Regla:} Ningún insumo o producto médico puede registrar un stock negativo. Cuando el inventario de un artículo alcance su umbral mínimo configurado, el sistema debe emitir una alerta obligatoria de reabastecimiento.

\textbf{Propósito:} Evitar quedarse sin insumos básicos a mitad de la jornada médica.

\subsection*{RN-03: Bloqueo de Modificaciones tras Cierre de Caja}
\textbf{Regla:} Una vez que el Administrador o Gerente ejecuta el arqueo de caja diario, la Recepcionista no podrá editar, anular ni agregar cobros o recibos pertenecientes a esa fecha finalizada sin autorización administrativa.

\textbf{Propósito:} Garantizar la integridad financiera y evitar alteraciones de cobros a destiempo.

\subsection*{RN-04: Obligatoriedad de Prescripción Médica}
\textbf{Regla:} Los productos medicinales de venta en clínica catalogados bajo receta solo podrán asociarse al cobro o dispensación si cuentan con la vinculación formal de la prescripción emitida por el médico tratante.

\textbf{Propósito:} Cumplir con la regulación médica y evitar la venta no autorizada de fármacos.

\subsection*{RN-05: Cierre Obligatorio de Cita con Declaración de Insumos}
\textbf{Regla:} Para marcar una cita médica de procedimiento, como curación de heridas o vendajes, en el estado ``Atendida'', el sistema debe exigir la declaración obligatoria del tipo y cantidad de materiales consumidos.

\textbf{Propósito:} Automatizar el descuento inmediato de inventario y evitar pérdidas o mermas sin registrar.
